% One-page CV -- COMPILERS track (LLVM, MLIR, polyhedral, codegen, runtimes).
% Content is a selection from full-siddharth-bhat.md; style lives in cvstyle.sty.
\documentclass[9pt,a4paper]{extarticle}
\usepackage{cvstyle}

\begin{document}
\cvheader

\section{Summary}
Compiler engineer and researcher (PhD, Cambridge) with a decade of performance-critical \cpp{}
inside production compilers and low-level runtimes.
\#3 contributor to \textbf{Polly}, LLVM's polyhedral loop optimizer (121 commits, 6000 LoC of
\cpp{} to LLVM, Fortran support and unified-memory CUDA codegen);
primary author of \textbf{Lean's LLVM backend}, lowering an IR through instruction selection and
code generation down to native machine code;
\#2 contributor to \textbf{Asterius}, a Haskell$\to$WebAssembly compiler whose runtime work was
merged into GHC.
Five talks at LLVM developer meetings on IR design, dominance, side-effect modelling and
optimization verification.
Comfortable reasoning about the machine underneath the language --- cache behaviour, memory
models, SIMD, assembly --- and optimizing hot code to that level.
Fluent in \cpp{}, Lean, Haskell, Rust, Rocq.

\section{Expertise}
\begin{records}
\item \textbf{Loop optimization \& HPC (\cpp{}, LLVM, CUDA):} \#3 contributor to Polly ---
GPU code generation for Fortran workloads, unified-memory support in the CUDA backend,
polyhedral scheduling; tiling patterns and stencil optimization in PolyMage, PLUTO and ISL
(SC 2017); LLVM GSoC mentor (2016, 2018).
\item \textbf{Backends \& runtimes:} primary author of Lean's LLVM backend; reimplemented the
GHC runtime on top of WebAssembly for Asterius (merged upstream into GHC); lock-free atomic
synchronization for real-time audio in the PPSSPP \cpp{} emulator; rewrote a raytracer in Haskell
to match hand-tuned \cpp{} (Haskell Exchange '20).
\item \textbf{IR design \& semantics (MLIR):} led \texttt{lean-mlir}, formal semantics for MLIR
with machine-checked peephole rewriting (ITP 2024); designed SSA-based IRs for functional
(CGO 2022) and quantum (CC 2021) compilation; verified translation validation for rewrites at
arbitrary bitwidth.
\end{records}

\section{Selected Publications}
\begin{pubs}
\pub{Verifying Peephole Rewriting in SSA Compiler IRs}{\me, Alex Keizer, Chris Hughes, Andres Goens, Tobias Grosser}{ITP 2024}
\pub{Lambda the Ultimate SSA: optimizing functional programs in SSA}{\me, Tobias Grosser}{CGO 2022}
\pub{Guided Equality Saturation}{Thomas Koehler, Andres Goens, \me, Tobias Grosser, Phil Trinder, Michel Steuwer}{POPL 2024}
\pub{QSSA: an SSA-based IR for Quantum Computing}{Anurudh Peduri, \me, Tobias Grosser}{CC 2021}
\pub{Optimizing Geometric Multigrid Computation using a DSL Approach}{Vinay Vasista, Kumudha KN, \me, Uday Bondhugula}{Supercomputing (SC) 2017}
\pub{Certified Decision Procedures for Width-Independent Bitvector Predicates}{\me, L\'eo Stefanesco, Chris Hughes, Tobias Grosser}{OOPSLA 2025}
\pub{Interactive Bit Vector Reasoning using Verified Bitblasting}{Henrik B\"oving, \me, Alex Keizer, Luisa Cicolini, Abdalrhman Mohamed, L\'eo Stefanesco, Josh Clune, Clark Barrett, Tobias Grosser}{OOPSLA 2025}
\end{pubs}
\vspace{1pt}
Full list (14 papers, incl.\ POPL, ICML, ICSE best paper, NeurIPS) on
\href{https://scholar.google.com/citations?user=7Irb-OUAAAAJ&hl=en}{Google Scholar}.

\section{Software}
\proj{Polly (LLVM)}{https://polly.llvm.org/}{\#3 contributor to LLVM's polyhedral loop optimizer: 121 commits, 6000 LoC, Fortran support and unified-memory CUDA codegen.}
\proj{Lean 4}{https://github.com/leanprover/lean4/pulls?q=author\%3Abollu}{primary author of the LLVM backend (instruction selection, codegen); top-20 contributor overall.}
\proj{Asterius / GHC}{https://github.com/tweag/asterius/commits?author=bollu}{\#2 contributor to this Haskell$\to$WebAssembly compiler; reimplemented the GHC runtime on WASM, later merged into GHC proper.}
\proj{lean-mlir}{https://github.com/opencompl/lean-mlir}{formal semantics for MLIR in Lean 4, with verified peephole rewriting.}
\proj{PLUTO \& PolyMage}{http://mcl.csa.iisc.ac.in/polymage.html}{loop-nest optimizers: found and fixed diamond-tiling bugs in PLUTO, added stencil and time-iterated-stencil support to PolyMage.}
\proj{Simplexhc}{http://github.com/bollu/simplexhc}{a compiler for a lazy functional language applying polyhedral ideas to an LLVM backend.}

\section{Experience}
\role{Sep--Nov '24}{Amazon Web Services, Automated Reasoning Group, Austin}{Research Intern}{Deciding memory non-interference in \texttt{lnsym}, a Lean-based ARM symbolic simulator; large-scale symbolic execution and its performance.}
\role{Jul--Sep '23}{Microsoft Research, Redmond}{Research Intern}{Retrieval-augmented proof synthesis for the \fstar{} proof assistant (ICSE 2025 best paper).}
\role{May--Jul '19}{Tweag.io, Paris}{Research Intern}{Re-implemented portions of the GHC runtime for \href{https://github.com/tweag/asterius/commits?author=bollu}{Asterius} (Haskell, C, WASM).}
\role{Summer '18}{Google Summer of Code, Polly Labs}{Mentor}{Enabling Polly's loop optimizations in Chapel.}
\role{Mar--Dec '17}{ETH Zurich, Scalable Parallel Computing}{Research Intern}{CUDA code generation in Polly for Fortran loop nests.}
\role{May--Jul '16}{IISc Bangalore}{Research Intern}{PolyMage DSL compiler; tiling patterns and stencil optimization; contributions to ISL and PLUTO.}

\section{Talks at LLVM Developer Meetings}
\href{https://www.youtube.com/watch?v=WtsInfbzxjs}{How to trust your peephole rewrites: automatically verifying them for arbitrary width} (EuroLLVM '25)
$\cdot$ \href{https://www.youtube.com/watch?v=4lh2NnLOxvQ}{lean-mlir: verifying peephole optimizations in MLIR} (LLVM Dev '24)
$\cdot$ \href{https://www.youtube.com/watch?v=VJORFvHJKWE}{(Correctly) extending dominance to MLIR regions} (LLVM Dev '23)
$\cdot$ \href{https://www.youtube.com/watch?v=6bDKasLZyxU}{MLIR side-effect modelling} (LLVM Dev '23)
$\cdot$ \href{https://www.youtube.com/watch?v=cyMQbZ0B84Q}{MLIR for functional programming} (EuroLLVM '22).\par

\section{Education}
\edu{PhD, Computer Science}{University of Cambridge; transferred from University of Edinburgh}{2022--2026 (writing up)}
\edu{MS, Computer Science}{IIIT Hyderabad, India}{2020--2021}
\edu{BTech, Computer Science}{IIIT Hyderabad, India}{2015--2020}

\end{document}
